% modules/3_proceso.tex
% ================================================================
% CAPÍTULO 3 - PROCESO DE MANTENIMIENTO
% Referencia: SWEBOK V4.0a s2 / ISO/IEC/IEEE 14764
% ================================================================

\chapter{Proceso de Mantenimiento}

La sección 2 del SWEBOK V4.0a describe el proceso de mantenimiento como un conjunto
de actividades estructuradas que garantizan que cada modificación al software preserve
su integridad. A continuación se define el proceso adoptado para el proyecto Yoliztli.

\section{Ciclo de Vida de una Solicitud de Modificación (MR)}

Toda solicitud de modificación --- ya sea una corrección de defecto (Problem Report --- PR)
o una mejora (Enhancement Request) --- sigue el ciclo de seis etapas definido en
ISO/IEC/IEEE 14764:

\begin{enumerate}
    \item \textbf{Registro de la solicitud:} El problema o mejora se documenta en el
          gestor de issues del repositorio Git con una descripción clara, el componente
          afectado y el tipo de mantenimiento (Correctivo, Preventivo, Adaptativo,
          Aditivo o Perfectivo).

    \item \textbf{Análisis de Impacto:} Antes de iniciar cualquier codificación, se
          evalúa cómo el cambio propuesto afecta a otros módulos del sistema, a la base
          de datos y a las pruebas existentes. Este análisis determina la estimación de
          esfuerzo.

    \item \textbf{Diseño e Implementación:} El desarrollo se realiza en una rama Git
          dedicada siguiendo la convención de nombres:
          \begin{itemize}
              \item \texttt{bugfix/<descripcion>} para mantenimiento correctivo.
              \item \texttt{feat/<descripcion>} para mantenimiento aditivo.
              \item \texttt{refactor/<descripcion>} para mantenimiento perfectivo.
              \item \texttt{chore/<descripcion>} para mantenimiento adaptativo y preventivo.
          \end{itemize}

    \item \textbf{Revisión y Aceptación:} Se ejecuta la suite de pruebas automatizadas
          (\texttt{php artisan test}) y se verifica manualmente el caso de prueba
          afectado. El código se fusiona a \texttt{main} solo si todas las pruebas pasan.

    \item \textbf{Migración y Despliegue:} Se aplican las nuevas migraciones de base
          de datos (\texttt{php artisan migrate}) en el entorno de destino y se
          recompila el frontend (\texttt{npm run build}).

    \item \textbf{Cierre y Documentación:} El issue se cierra en el gestor de
          repositorio y se documenta el cambio en el historial de revisiones de este
          plan.
\end{enumerate}

\newpage
\section{Convención de Commits}

Todos los commits del proyecto siguen la especificación de \textbf{Conventional Commits}
para mantener un historial de cambios legible y estructurado:

\begin{table}[H]
\centering
\begin{tabular}{lll}
\toprule
\textbf{Prefijo} & \textbf{Categoría SWEBOK} & \textbf{Ejemplo de uso} \\
\midrule
\texttt{fix:}      & Correctivo  & \texttt{fix: corregir almacenamiento en tabla progress} \\
\texttt{feat:}     & Aditivo     & \texttt{feat: agregar filtros por edad en dashboard} \\
\texttt{refactor:} & Perfectivo  & \texttt{refactor: optimizar consultas en DocenteController} \\
\texttt{chore:}    & Adaptativo / Preventivo & \texttt{chore: migrar conexion a postgresql} \\
\texttt{docs:}     & Documentación & \texttt{docs: actualizar plan de mantenimiento} \\
\texttt{test:}     & Preventivo  & \texttt{test: agregar pruebas para ProgressController} \\
\bottomrule
\end{tabular}
\caption{Convención de mensajes de commit por categoría de mantenimiento}
\label{tab:commits}
\end{table}

\newpage
\section{Gestión de la Configuración del Software (SCM)}

De acuerdo con la relación que SWEBOK establece entre mantenimiento y
\textit{Software Configuration Management} (SCM), el repositorio Git del proyecto
aplica las siguientes prácticas:

\begin{itemize}
    \item \textbf{Control de versiones centralizado:} Repositorio alojado en GitHub
          bajo la organización \texttt{E33-a/Estancia-I}.
    \item \textbf{Rama estable:} La rama \texttt{main} representa siempre el estado
          del sistema en producción. Ningún cambio se fusiona a esta rama sin haber
          pasado las pruebas de regresión.
    \item \textbf{Archivo \texttt{.gitignore}:} Excluye credenciales sensibles
          (\texttt{.env}), dependencias de terceros (\texttt{/vendor},
          \texttt{/node\_modules}) y archivos de base de datos locales
          (\texttt{*.sqlite}).
    \item \textbf{Etiquetado de versiones:} Cada entrega al catedrático se etiqueta
          con \texttt{git tag} siguiendo el esquema \texttt{v1.x.y} (versionado
          semántico).
\end{itemize}
